\section{Introduction}
\label{sec:introduction}

Large Language Models (LLMs) have demonstrated remarkable capabilities in natural language understanding and generation~\citep{brown2020language}, including complex multi-step reasoning when prompted with chain-of-thought strategies~\citep{wei2022chain}. These models acquire extensive world knowledge during pre-training, storing factual associations within their parameters---commonly referred to as \emph{parametric knowledge}. However, parametric knowledge is inherently static: it reflects the training data distribution and cannot be updated without retraining.

Retrieval-Augmented Generation (RAG)~\citep{lewis2020retrieval} addresses this limitation by supplementing LLMs with external documents retrieved at inference time. While RAG generally improves factual accuracy~\citep{guu2020realm}, a critical failure mode arises when retrieved context \emph{contradicts} the model's parametric knowledge---a \emph{knowledge conflict}. Such conflicts occur naturally when retrieved documents contain outdated information, factual errors, or adversarially manipulated content~\citep{pan2023risk}. When faced with a conflict, the model must implicitly decide whether to trust its parametric memory or defer to the provided context, with significant implications for system reliability.

This problem becomes substantially more complex in \emph{multi-hop reasoning}, where answering a question requires synthesizing information across multiple documents in a reasoning chain~\citep{yang2018hotpotqa, trivedi2022musique}. A conflict at any single hop can potentially propagate through the chain, corrupting downstream inferences. Despite growing interest in knowledge conflicts~\citep{longpre2021entity, xie2024adaptive, chen2022rich}, existing work has focused primarily on single-hop settings. The interaction between conflicts and hop position in multi-step reasoning chains remains poorly understood.

This work addresses this gap through a systematic experimental study. We inject controlled conflicts at specific hop positions and evaluate four LLMs across two multi-hop datasets, investigating six research questions:

\begin{enumerate}
    \item[\textbf{RQ1}] How significantly do knowledge conflicts degrade multi-hop QA accuracy?
    \item[\textbf{RQ2}] Does the hop position of conflict injection affect degradation in 2-hop reasoning?
    \item[\textbf{RQ3}] How does error propagation differ between 2-hop and 3-hop chains?
    \item[\textbf{RQ4}] Does model size predict robustness to conflicts?
    \item[\textbf{RQ5}] Does the conflict type (factual, temporal, numerical) influence severity?
    \item[\textbf{RQ6}] Does entity popularity affect conflict resolution strategy?
\end{enumerate}

\noindent\textbf{Contributions.} We make five contributions: (1)~a controlled conflict injection framework for multi-hop QA; (2)~a large-scale evaluation across 4 models, 2 datasets, and 3 conflict types totaling 19,416 examples; (3)~the discovery that answer-bearing hop conflicts cause 40--52pp drops while bridge-hop conflicts are far less damaging---a pattern consistent across 2-hop and 3-hop chains; (4)~a conflict type taxonomy showing temporal conflicts are most disruptive (50.8pp drop, CFR up to 80\%); and (5)~evidence that entity popularity influences conflict resolution.
